Modularity is a key quality attribute in software systems, being directly related to maintainability, understandability, and evolution. In this context, architectural decisions play a fundamental role in the structural organization of systems, influencing aspects such as coupling, cohesion, and dependencies among components. However, the selection of architectural styles, such as layered architecture and hexagonal architecture, is often based on best practices and prior experience, without a consistent quantitative evaluation of their impact on software structural quality.

Therefore, this work aims to evaluate the modularity of object-oriented systems through the application of structural metrics across different architectural styles. To achieve this, a comparative experimental study is conducted using two previously developed systems, each representing a distinct architectural approach. The analysis is based on classical structural metrics related to coupling, cohesion, and complexity, enabling a quantitative assessment of the systems' structural quality.

This work is expected to contribute to a better understanding of how architectural decisions affect software modularity, providing support for more evidence-based practices in Software Engineering.

% Separe as Keywords por ponto e vírgula.
\keywords{ Modularity; Software Architecture; Software Metrics; Object-Oriented Systems.}